% GENERATED FILE. Edit canonical lessons, then run npm run generate:books.
% canonical-source-hash: fnv1a64:4327a0f2955fe43f
% canonical-lessons: TE-C524-pratibha, TE-C524-pravinyam, TE-C524-panitanam, TE-C524-nerpari, TE-C524-corava

\chapter{Skill and Talent}
\label{ch:te-skill-and-talent}
\hlchaptermodality{524}

\begin{chapteropening}
\textbf{By the end of this chapter:} \emph{I can talk about talent, mastery, craft and initiative.}

Complete the last lesson of chapter 524: I can talk about talent, mastery, craft and initiative.
\end{chapteropening}

\section[\texorpdfstring{\te{ప్రతిభ} (pratibha)}{ప్రతిభ (pratibha)}]{\te{ప్రతిభ} (pratibha) — talent; flair}

\label{lesson:TE-C524-pratibha}



\begin{quote}
Before the new one: say the Telugu for an assistant, then the Telugu for a contractor.
\end{quote}

\subsection*{You'll want to know: \te{ప్రతిభ}}
\textbf{\te{ప్రతిభ}} — \emph{pratibha} — \textquotedblleft{}talent; flair\textquotedblright{}.

\begin{quote}
\te{ఆ} \te{పిల్లవాడికి} \te{సంగీతంలో} \te{గొప్ప} \te{ప్రతిభ} \te{ఉంది}.
\end{quote}

\begin{quote}
\emph{ā pillavāḍiki saṅgītaṁlō goppa pratibha undi.}
\end{quote}

\begin{quote}
That boy has a great talent for music.
\end{quote}

\te{ప్రతిభావంతుడు}, a gifted man.

\begin{cousinweb}
\textbf{\te{ప్రతిభ}} (\emph{pratibha}) is borrowed from Sanskrit \emph{pratibhā}, \textquotedblleft{}light; intelligence, understanding\textquotedblright{} (Monier-Williams).
\end{cousinweb}

\begin{grammarlens}[title={What you've built}]
One more word for describing what you lived and felt.
\end{grammarlens}

\subsection*{Your turn}
\begin{itemize}
\raggedright
  \item \emph{Say it:} \emph{pratibha}
  \item \emph{Say it:} \emph{pratibha}, once more
  \item \emph{Say it:} \emph{ā pillavāḍiki saṅgītaṁlō goppa pratibha undi.} — That boy has a great talent for music.
  \item \emph{Recall:} say the Telugu for an assistant, then the Telugu for a contractor, then say \emph{pratibha} again
\end{itemize}

\begin{tcolorbox}[breakable,colback=teal!4,colframe=teal!35!black,title={Before you move on}]
What does \te{ప్రతిభ} mean? (\textquotedblleft{}Talent; flair\textquotedblright{}.) Say it once more.
\end{tcolorbox}
\section[\texorpdfstring{\te{ప్రావీణ్యం} (prāvīṇyaṁ)}{ప్రావీణ్యం (prāvīṇyaṁ)}]{\te{ప్రావీణ్యం} (prāvīṇyaṁ) — proficiency; mastery}

\label{lesson:TE-C524-pravinyam}



\begin{quote}
Before the new one: say the Telugu for a contractor, then the Telugu for talent.
\end{quote}

\subsection*{You'll want to know: \te{ప్రావీణ్యం}}
\textbf{\te{ప్రావీణ్యం}} — \emph{prāvīṇyaṁ} — \textquotedblleft{}proficiency; mastery\textquotedblright{}.

\begin{quote}
\te{ఆమెకు} \te{మూడు} \te{భాషల్లో} \te{ప్రావీణ్యం} \te{ఉంది}.
\end{quote}

\begin{quote}
\emph{āmeku mūḍu bhāṣallō prāvīṇyaṁ undi.}
\end{quote}

\begin{quote}
She is proficient in three languages.
\end{quote}

-\te{లో} \te{ప్రావీణ్యం}, skill in.

\begin{cousinweb}
\textbf{\te{ప్రావీణ్యం}} (\emph{prāvīṇyaṁ}) is borrowed from Sanskrit \emph{prāvīṇya}, \textquotedblleft{}cleverness, dexterity, proficiency\textquotedblright{} (Monier-Williams).
\end{cousinweb}

\begin{grammarlens}[title={What you've built}]
One more word for describing what you lived and felt.
\end{grammarlens}

\subsection*{Your turn}
\begin{itemize}
\raggedright
  \item \emph{Say it:} \emph{prāvīṇyaṁ}
  \item \emph{Say it:} \emph{prāvīṇyaṁ}, once more
  \item \emph{Say it:} \emph{āmeku mūḍu bhāṣallō prāvīṇyaṁ undi.} — She is proficient in three languages.
  \item \emph{Recall:} say the Telugu for a contractor, then the Telugu for talent, then say \emph{prāvīṇyaṁ} again
\end{itemize}

\begin{tcolorbox}[breakable,colback=teal!4,colframe=teal!35!black,title={Before you move on}]
What does \te{ప్రావీణ్యం} mean? (\textquotedblleft{}Proficiency; mastery\textquotedblright{}.) Say it once more.
\end{tcolorbox}
\section[\texorpdfstring{\te{పనితనం} (panitanaṁ)}{పనితనం (panitanaṁ)}]{\te{పనితనం} (panitanaṁ) — workmanship; craft}

\label{lesson:TE-C524-panitanam}



\begin{quote}
Before the new one: say the Telugu for talent, then the Telugu for proficiency.
\end{quote}

\subsection*{You'll want to know: \te{పనితనం}}
\textbf{\te{పనితనం}} — \emph{panitanaṁ} — \textquotedblleft{}workmanship; craft\textquotedblright{}.

\begin{quote}
\te{ఆ} \te{కుర్చీలో} \te{వడ్రంగి} \te{పనితనం} \te{కనిపిస్తుంది}.
\end{quote}

\begin{quote}
\emph{ā kurcīlō vaḍraṅgi panitanaṁ kanipistundi.}
\end{quote}

\begin{quote}
The carpenter's craft shows in that chair.
\end{quote}

\te{పని}, work, + -\te{తనం}.

\begin{cousinweb}
\textbf{\te{పనితనం}} (\emph{panitanaṁ}) is native Dravidian: DEDR 3884 lists Telugu \emph{pani}, \textquotedblleft{}work, labour, workmanship, art\textquotedblright{}, with Tamil \emph{paṇi}; -\te{తనం} makes the abstract noun.
\end{cousinweb}

\begin{grammarlens}[title={What you've built}]
One more word for describing what you lived and felt.
\end{grammarlens}

\subsection*{Your turn}
\begin{itemize}
\raggedright
  \item \emph{Say it:} \emph{panitanaṁ}
  \item \emph{Say it:} \emph{panitanaṁ}, once more
  \item \emph{Say it:} \emph{ā kurcīlō vaḍraṅgi panitanaṁ kanipistundi.} — The carpenter's craft shows in that chair.
  \item \emph{Recall:} say the Telugu for talent, then the Telugu for proficiency, then say \emph{panitanaṁ} again
\end{itemize}

\begin{tcolorbox}[breakable,colback=teal!4,colframe=teal!35!black,title={Before you move on}]
What does \te{పనితనం} mean? (\textquotedblleft{}Workmanship; craft\textquotedblright{}.) Say it once more.
\end{tcolorbox}
\section[\texorpdfstring{\te{నేర్పరి} (nērpari)}{నేర్పరి (nērpari)}]{\te{నేర్పరి} (nērpari) — a skilled person; an adept}

\label{lesson:TE-C524-nerpari}



\begin{quote}
Before the new one: say the Telugu for proficiency, then the Telugu for workmanship.
\end{quote}

\subsection*{You'll want to know: \te{నేర్పరి}}
\textbf{\te{నేర్పరి}} — \emph{nērpari} — \textquotedblleft{}a skilled person; an adept\textquotedblright{}.

\begin{quote}
\te{మా} \te{అమ్మ} \te{వంటలో} \te{నేర్పరి}.
\end{quote}

\begin{quote}
\emph{mā amma vaṇṭalō nērpari.}
\end{quote}

\begin{quote}
My mother is an expert cook.
\end{quote}

\te{నేర్పు}, skill, + -\te{అరి}, one who has it.

\begin{cousinweb}
\textbf{\te{నేర్పరి}} (\emph{nērpari}) is native Dravidian: DEDR 3770 lists Telugu \emph{nēr(u)pu}, \textquotedblleft{}cleverness, skill\textquotedblright{}, beside \emph{nēr(u)cu}, \textquotedblleft{}to learn\textquotedblright{}; the ending -\te{అరి} makes \textquotedblleft{}one who has it\textquotedblright{}.
\end{cousinweb}

\begin{grammarlens}[title={What you've built}]
One more word for describing what you lived and felt.
\end{grammarlens}

\subsection*{Your turn}
\begin{itemize}
\raggedright
  \item \emph{Say it:} \emph{nērpari}
  \item \emph{Say it:} \emph{nērpari}, once more
  \item \emph{Say it:} \emph{mā amma vaṇṭalō nērpari.} — My mother is an expert cook.
  \item \emph{Recall:} say the Telugu for proficiency, then the Telugu for workmanship, then say \emph{nērpari} again
\end{itemize}

\begin{tcolorbox}[breakable,colback=teal!4,colframe=teal!35!black,title={Before you move on}]
What does \te{నేర్పరి} mean? (\textquotedblleft{}A skilled person; an adept\textquotedblright{}.) Say it once more.
\end{tcolorbox}
\section[\texorpdfstring{\te{చొరవ} (corava)}{చొరవ (corava)}]{\te{చొరవ} (corava) — initiative; drive}

\label{lesson:TE-C524-corava}



\begin{quote}
Before the new one: say the Telugu for workmanship, then the Telugu for a skilled person.
\end{quote}

\subsection*{You'll want to know: \te{చొరవ}}
\textbf{\te{చొరవ}} — \emph{corava} — \textquotedblleft{}initiative; drive\textquotedblright{}.

\begin{quote}
\te{ఆమె} \te{చొరవ} \te{తీసుకుని} \te{సమస్యను} \te{పరిష్కరించింది}.
\end{quote}

\begin{quote}
\emph{āme corava tīsukuni samasyanu pariṣkariñcindi.}
\end{quote}

\begin{quote}
She took the initiative and solved the problem.
\end{quote}

\te{చొరవ} \te{తీసుకో}, to take the initiative.

\begin{cousinweb}
\textbf{\te{చొరవ}} (\emph{corava}) is a Telugu word for the drive to step forward and act; its origin is not settled.
\end{cousinweb}

\begin{grammarlens}[title={What you've built}]
One more word for describing what you lived and felt.
\end{grammarlens}

\subsection*{Your turn}
\begin{itemize}
\raggedright
  \item \emph{Say it:} \emph{corava}
  \item \emph{Say it:} \emph{corava}, once more
  \item \emph{Say it:} \emph{āme corava tīsukuni samasyanu pariṣkariñcindi.} — She took the initiative and solved the problem.
  \item \emph{Recall:} say the Telugu for workmanship, then the Telugu for a skilled person, then say \emph{corava} again
\end{itemize}

\begin{tcolorbox}[breakable,colback=teal!4,colframe=teal!35!black,title={Before you move on}]
What does \te{చొరవ} mean? (\textquotedblleft{}Initiative; drive\textquotedblright{}.) Say it once more.
\end{tcolorbox}
